\section{Architektury Modeli}
\label{sec:architectures}

W projekcie zaimplementowano i porównano trzy różne architektury sieci neuronowych, każda z unikalnymi cechami dostosowanymi do pracy z danymi tabelarycznymi.

\subsection{SimpleMLP}

Klasyczna architektura fully-connected: warstwa wejściowa (22 cechy) → warstwa ukryta → dropout → warstwa wyjściowa.

\begin{table}[h]
\centering
\caption{Przestrzeń hiperparametrów dla MLP}
\label{tab:hyperparams_mlp}
\begin{tabular}{lll}
\toprule
Hiperparametr & Wartości & Opis \\
\midrule
hidden\_channels & [64, 128, 256] & Liczba neuronów w warstwie ukrytej \\
activation\_layer & ['relu', 'tanh'] & Funkcja aktywacji \\
dropout & [0.2, 0.3, 0.4] & Współczynnik dropout \\
\bottomrule
\end{tabular}
\end{table}


\subsection{TabNet}

Specjalizowana architektura wykorzystująca mechanizm attention do sekwencyjnej selekcji cech \cite{arik2019tabnet}. Oferuje interpretowalność poprzez maski attention pokazujące istotność cech.

\begin{table}[h]
\centering
\caption{Przestrzeń hiperparametrów dla TABNET}
\label{tab:hyperparams_tabnet}
\begin{tabular}{lll}
\toprule
Hiperparametr & Wartości & Opis \\
\midrule
n\_d & [8, 32, 64] & Wymiar feature transformers \\
n\_a & [8, 32, 64] & Wymiar attention layers \\
n\_steps & [3, 5, 7] & Liczba kroków decyzyjnych \\
\bottomrule
\end{tabular}
\end{table}


\subsection{TabTransformer}

Architektura oparta na mechanizmie Transformer \cite{vaswani2017attention}, wykorzystująca self-attention do modelowania zależności między cechami.

\begin{table}[h]
\centering
\caption{Przestrzeń hiperparametrów dla TABTRANSFORMER}
\label{tab:hyperparams_tabtransformer}
\begin{tabular}{lll}
\toprule
Hiperparametr & Wartości & Opis \\
\midrule
d\_model & [64, 128, 256] & Wymiar modelu \\
n\_heads & [2, 4, 8] & Liczba głów attention \\
num\_layers & [1, 2, 3] & Liczba warstw transformera \\
\bottomrule
\end{tabular}
\end{table}

